\documentclass[12pt,a4paper]{article}

\usepackage[utf8]{inputenc}
\usepackage[T1]{fontenc}
\usepackage[english]{babel}
\usepackage{amsmath,amssymb,bm}
\usepackage{graphicx}
\usepackage{float}
\usepackage{booktabs}
\usepackage{geometry}
\usepackage[expansion=false]{microtype}
\usepackage{hyperref}
\usepackage{xcolor}

\geometry{margin=2.2cm}
\setlength{\emergencystretch}{2em}

% Figures are kept at the place where they are introduced in the text.
\hypersetup{colorlinks=true,linkcolor=black,urlcolor=blue}

\graphicspath{{figures/}{./}}

\newcommand{\reportfigure}[3]{%
\begin{figure}[H]
    \centering
    \includegraphics[width=0.94\textwidth]{#1}%
    \caption{#2}
    \label{#3}
\end{figure}
}

\title{Dispersion of Exciton--Photon Polaritons\\ in a Planar DBR Microcavity}
\author{Korostelev Nikita}
\date{}

\begin{document}
\maketitle
\clearpage

\section{Aim of the work}

The aim of this work is to calculate the dispersion of exciton--photon polaritons in a planar semiconductor microcavity and to understand how it arises from the coupling between light in the cavity and an exciton in a quantum well. Analytical models are used to understand the origin of the main dependences, and then they are compared with the calculation of the full multilayer structure by the transfer-matrix method.

For oblique propagation, TE and TM polarizations are considered separately. For TE polarization, the electric field is perpendicular to the plane of incidence. For TM polarization, the magnetic field is perpendicular to this plane.

The main physical system consists of the photon mode of the microcavity and the exciton in the quantum well. In the weak-coupling regime, their energies are only weakly affected by each other. In the strong-coupling regime, new mixed states appear: the lower and upper polaritons.

\clearpage
\section{Microcavity structure and parameters}

The following structure is considered:

\[
\mathrm{air}\,|\,(HL)^{23}\,|\,C_1\,|\,W\,|\,C_2\,|\,(LH)^{26}\,|\,\mathrm{air},
\]

where $H$ is the high-refractive-index GaAs layer, $L$ is the lower-index AlAs layer, $C_1$ and $C_2$ are the two parts of the GaAs cavity layer, and $W$ is the In$_{0.08}$Ga$_{0.92}$As quantum well. Light is incident from the left through the top DBR. On both sides of the cavity, an $L$ layer is adjacent to the GaAs cavity layer.

\reportfigure{01_refractive_index_profile.png}{Refractive-index profile $n(z)$ of the multilayer microcavity. The figure is built directly from the chosen geometry and refractive indices, so it is not a result of solving the dispersion problem. The red line marks the center of the QW.}{fig:index-profile}

Figure~\ref{fig:index-profile} shows two finite DBR mirrors and the wider cavity region between them. The QW is very thin compared with the total length of the structure, so on this scale it mainly indicates its position inside the cavity.

As an experimental reference, we use the structure reported by P. Cilibrizzi et al., \textit{Polariton condensation in a strain-compensated planar microcavity with InGaAs quantum wells}, Applied Physics Letters 105, 191118 (2014): a $2\lambda$ cavity, 23 and 26 DBR periods, $6$~nm quantum wells, and a total Rabi splitting of about $8$~meV for six strongly coupled QWs.

For the structure considered here, constant refractive indices are used:

\[
n_H=3.60,\qquad n_L=2.95,\qquad n_C=3.60,\qquad n_W=3.65.
\]

The exciton resonance energy at zero in-plane wave vector is fixed as

\[
E_X(0)=1.450\ \mathrm{eV}.
\]

Here $E_X(0)$ is the exciton energy at $k_\parallel=0$. The value $1.450$~eV is in the same energy range, around $1.45$~eV, where Cilibrizzi et al. observed polariton branches. It sets the center of the spectral range considered here and is not a result of our calculation.

The corresponding wavelength in vacuum is

\[
\lambda_0=\frac{hc}{E_X(0)}\approx0.855\ \mu\mathrm m.
\]

The initial thicknesses of the quarter-wave DBR layers are

\[
d_H=\frac{\lambda_0}{4n_H}\approx59.4\ \mathrm{nm},
\qquad
 d_L=\frac{\lambda_0}{4n_L}\approx72.5\ \mathrm{nm}.
\]

For a $2\lambda$ cavity, the initial estimate of the total thickness of the cavity layer is

\[
L_C^{(0)}=\frac{2\lambda_0}{n_C}\approx475\ \mathrm{nm}.
\]

The final value of $L_C$ is determined from the resonance of the full DBR structure. The reason is that reflection from a DBR adds an extra phase shift and slightly changes the resonance condition.

The quantum well has thickness

\[
d_{QW}=6\ \mathrm{nm}
\]

and replaces a GaAs region of the same thickness inside the cavity. Therefore,

\[
L_C=d_{C_1}+d_{QW}+d_{C_2}.
\]

\clearpage
\section{Dispersion of uncoupled modes}

\subsection{Cavity photon mode}

In a planar cavity, the longitudinal component of the wave vector is fixed by the resonance condition, while $k_\parallel$ remains continuous. Therefore, the photon energy is

\[
E_C(k_\parallel)
=
\frac{\hbar c}{n_C}
\sqrt{k_z^2+k_\parallel^2}.
\]

Using the mode energy at normal propagation, this expression can be written as

\[
E_C(k_\parallel)
=
\sqrt{
E_C(0)^2+
\left(\frac{\hbar c}{n_C}k_\parallel\right)^2
}.
\]

For small $k_\parallel$, the dispersion becomes parabolic:

\[
E_C(k_\parallel)
\approx
E_C(0)+\frac{\hbar^2k_\parallel^2}{2m_{ph}},
\]

where

\[
m_{ph}=\frac{E_C(0)n_C^2}{c^2}.
\]

The parabolic form does not come from a nonzero rest mass of a free photon. It appears because the cavity fixes $k_z$: when only the in-plane wave vector changes, the square-root expression can be expanded in powers of $k_\parallel^2$.

\subsection{Exciton in a quantum well}

The quantum well confines the motion of the electron and hole along $z$, but the exciton center of mass remains free in the plane of the layer. Therefore, its dispersion is

\[
E_X(k_\parallel)
=
E_X(0)+\frac{\hbar^2k_\parallel^2}{2M_X},
\]

where $M_X$ is the exciton center-of-mass mass. The calculation uses

\[
M_X=0.5m_0.
\]

Since $M_X\gg m_{ph}$, the exciton branch is much flatter than the photon branch. The analytical plots use the same range $k_\parallel=0\ldots4\ \mu\mathrm m^{-1}$ as the later numerical calculations.

These two uncoupled dispersions are not shown in a separate figure. They are plotted together with the polariton branches in the next figure. This makes it possible to see directly how the crossing of the bare photon and exciton modes becomes an avoided crossing after the coupling is switched on.

\clearpage
\section{Analytical polariton model}

For a fixed in-plane wave vector $k_\parallel$, two degrees of freedom are considered: a photon in the microcavity and an exciton in the quantum well. In Kavokin et al., \textit{Microcavities}, Sections~5.8.1--5.8.2, this system is written as two coupled quantum modes:

\[
\hat H
=
\sum_{k_\parallel}
E_C(k_\parallel)\hat B_{k_\parallel}^{\dagger}\hat B_{k_\parallel}
+
\sum_{k_\parallel}
E_X(k_\parallel)\hat X_{k_\parallel}^{\dagger}\hat X_{k_\parallel}
+
\sum_{k_\parallel}
g(k_\parallel)
\left(
\hat B_{k_\parallel}^{\dagger}\hat X_{k_\parallel}
+
\hat X_{k_\parallel}^{\dagger}\hat B_{k_\parallel}
\right).
\]

Here $\hat B_{k_\parallel}$ annihilates one photon in the cavity mode, while $\hat B_{k_\parallel}^{\dagger}$ creates one. Similarly, $\hat X_{k_\parallel}$ and $\hat X_{k_\parallel}^{\dagger}$ annihilate and create an exciton. The quantities $E_C(k_\parallel)$ and $E_X(k_\parallel)$ are the uncoupled photon and exciton energies, and $g$ is their coupling energy. The same $k_\parallel$ is conserved for the photon and exciton because the structure is uniform in the plane of the layers.

In the considered range, the coupling is taken to be constant:

\[
g(k_\parallel)\approx g,
\qquad
\Omega_R=2g,
\]

where $\Omega_R$ is the Rabi splitting: the energy separation between the upper and lower polariton branches at the point where the bare photon and exciton modes would have the same energy.

For each fixed $k_\parallel$, the Hamiltonian reduces to the matrix

\[
H(k_\parallel)=
\begin{pmatrix}
E_C(k_\parallel) & g\\
g & E_X(k_\parallel)
\end{pmatrix}
=
\begin{pmatrix}
E_C(k_\parallel) & \Omega_R/2\\
\Omega_R/2 & E_X(k_\parallel)
\end{pmatrix}.
\]

The detuning between the uncoupled modes at an arbitrary $k_\parallel$ is defined as

\[
\Delta(k_\parallel)
=
E_C(k_\parallel)-E_X(k_\parallel).
\]

It is also useful to define the detuning at normal propagation:

\[
\delta
=
E_C(0)-E_X(0).
\]

If both uncoupled dispersions are written in the parabolic approximation, then

\[
\Delta(k_\parallel)
=
\delta
+
\frac{\hbar^2k_\parallel^2}{2}
\left(
\frac{1}{m_{ph}}-\frac{1}{M_X}
\right),
\]

where $m_{ph}$ is the effective cavity-photon mass and $M_X$ is the exciton center-of-mass mass in the QW plane. Since the exciton mass $M_X$ is much larger than the effective photon mass $m_{ph}$, the exciton branch changes slowly. Therefore, the change of detuning with $k_\parallel$ is mainly set by the photon branch.

The eigenenergies are found from

\[
\det\!\left[H(k_\parallel)-E I\right]=0,
\]

which gives

\[
\left[E_C(k_\parallel)-E\right]
\left[E_X(k_\parallel)-E\right]
-g^2=0.
\]

The two solutions form the upper and lower polariton branches:

\[
E_{UP,LP}(k_\parallel)
=
\frac{E_C(k_\parallel)+E_X(k_\parallel)}{2}
\pm
\frac12
\sqrt{\Delta(k_\parallel)^2+\Omega_R^2}.
\]

The index $UP$ corresponds to the ``$+$'' sign and denotes the upper polariton branch. The index $LP$ corresponds to the ``$-$'' sign and denotes the lower polariton branch. At exact resonance, $\Delta=0$,

\[
E_{UP}-E_{LP}=\Omega_R.
\]

Therefore, the crossing of the uncoupled modes is replaced by an avoided crossing, or anticrossing.

\reportfigure{02_analytical_polariton_dispersion.png}{Analytical dispersion of the photon, exciton, and polariton branches at zero detuning. All four curves are obtained analytically: $E_C(k_\parallel)$ and $E_X(k_\parallel)$ describe the uncoupled modes, while $E_{LP}(k_\parallel)$ and $E_{UP}(k_\parallel)$ are the eigenvalues of the $2\times2$ matrix. The vertical axis shows energy relative to $E_X(0)$.}{fig:analytical-dispersion}

Figure~\ref{fig:analytical-dispersion} shows that the coupling replaces the crossing of the bare modes by an anticrossing. As $k_\parallel$ increases, the lower branch approaches the almost flat exciton dispersion, while the upper branch follows the photon dispersion more closely.

A polariton is a linear combination of photon and exciton states. In the notation of Section~5.8.1 of the book, the Hopfield transformation is

\[
\hat a_{LP,k}
=
X_k\hat X_k-C_k\hat B_k,
\qquad
\hat a_{UP,k}
=
C_k\hat X_k+X_k\hat B_k.
\]

Here $\hat a_{LP,k}$ and $\hat a_{UP,k}$ correspond to the lower and upper polaritons. The quantities $C_k$ and $X_k$ are called the Hopfield coefficients. They satisfy the normalization condition

\[
|C_k|^2+|X_k|^2=1.
\]

The quantity $|C|^2$ gives the photon fraction of the polariton, and $|X|^2$ gives the exciton fraction. For the lower branch,

\[
|C_{LP}(k_\parallel)|^2
=
\frac12
\left(
1-
\frac{\Delta(k_\parallel)}
{\sqrt{\Delta(k_\parallel)^2+\Omega_R^2}}
\right),
\]

\[
|X_{LP}(k_\parallel)|^2
=
\frac12
\left(
1+
\frac{\Delta(k_\parallel)}
{\sqrt{\Delta(k_\parallel)^2+\Omega_R^2}}
\right).
\]

For the upper branch, the fractions are interchanged:

\[
|C_{UP}(k_\parallel)|^2=|X_{LP}(k_\parallel)|^2,
\qquad
|X_{UP}(k_\parallel)|^2=|C_{LP}(k_\parallel)|^2.
\]

At $\Delta=0$, both polariton modes have equal photon and exciton fractions,

\[
|C|^2=|X|^2=\frac12.
\]

At large detuning, one branch becomes mainly photon-like, while the other becomes mainly exciton-like.

\reportfigure{03_hopfield_coefficients.png}{Photon and exciton fractions of the lower and upper polaritons. The curves are calculated analytically from the eigenvectors of the same $2\times2$ model; $|C|^2$ and $|X|^2$ give the photon and exciton fractions, respectively.}{fig:hopfield}

Figure~\ref{fig:hopfield} gives a simple way to read the later results: as $k_\parallel$ increases, the lower polariton becomes mainly exciton-like, while the upper polariton becomes mainly photon-like. This change of composition will later appear both in reflectivity and in the linewidths.

\clearpage
\section{Finite DBR microcavity}

The analytical parabola does not include the finite number of mirror layers, the phase shift on reflection, or the difference between TE and TM polarization at nonzero $k_\parallel$. Therefore, we next use the transfer-matrix method (TMM). Its idea is simple: propagation inside each layer and reflection at each interface are described separately, and then the matrices of all layers are multiplied. The general method is described in Kavokin et al., \textit{Microcavities}, Sections~2.5--2.7.

For layer $j$ with refractive index $n_j$ and thickness $d_j$, the longitudinal component of the wave vector is

\[
q_j(E,k_\parallel)
=
\sqrt{
\left(\frac{n_jE}{\hbar c}\right)^2-k_\parallel^2
}.
\]

When the wave propagates through a layer, its phase changes. For a layer of thickness $d_j$, the phase shift is

\[
\varphi_j=q_jd_j.
\]

Here $\varphi_j$ shows how many radians of phase are accumulated across layer $j$.

To write the interface conditions for TE and TM polarization in the same form, auxiliary quantities $Y_j$ are introduced:

\[
Y_j^{TE}=n_j\cos\phi_j,
\qquad
Y_j^{TM}=\frac{\cos\phi_j}{n_j},
\]

where $\phi_j$ is the propagation angle inside layer $j$, and

\[
\cos\phi_j=
\frac{q_j}{n_jE/(\hbar c)}.
\]

The quantity $Y_j$ is not a new material property. It is a convenient combination of $n_j$ and the angle $\phi_j$ that allows the boundary conditions for TE and TM waves to be written in the same form. These relations follow the notation of Section~2.5 of \textit{Microcavities}.

In the basis of forward and backward waves,

\[
\mathbf v(z)=
\begin{pmatrix}
A^+(z)\\
A^-(z)
\end{pmatrix},
\]

the propagation matrix through a homogeneous layer is

\[
P_j=
\begin{pmatrix}
e^{iq_jd_j} & 0\\
0 & e^{-iq_jd_j}
\end{pmatrix}.
\]

The interface matrix for $i\rightarrow j$ is

\[
I_{i\rightarrow j}
=
\frac{1}{2Y_j}
\begin{pmatrix}
Y_i+Y_j & Y_j-Y_i\\
Y_j-Y_i & Y_i+Y_j
\end{pmatrix}.
\]

The total matrix $M$ connects the amplitudes on the left and right sides:

\[
\mathbf v_R=M\mathbf v_L.
\]

For light incident from the left,

\[
\mathbf v_L=
\begin{pmatrix}
1\\
r
\end{pmatrix},
\qquad
\mathbf v_R=
\begin{pmatrix}
t\\
0
\end{pmatrix}.
\]

Therefore,

\[
r=-\frac{M_{21}}{M_{22}},
\qquad
R=|r|^2.
\]

First, the exciton response is switched off, while the InGaAs layer itself remains inside the cavity. This defines the photon resonance of exactly the same geometry without exciton--photon coupling. The cavity thickness $L_C$ is adjusted so that at $k_\parallel=0$,

\[
E_C^{bg}(0)=E_X(0)
\]

for zero detuning. This removes an additional shift of the photon mode caused by the difference between $n_W$ and $n_C$.

For the dispersion calculations, the range $k_\parallel=0\ldots4\ \mu\mathrm m^{-1}$ is used with a step of about $0.0167\ \mu\mathrm m^{-1}$. The original reflectivity spectrum is calculated on an energy grid with a step of $0.0005$~meV. To make very narrow resonances visible on the map, the displayed value is the minimum reflectivity inside an energy bin of width $0.050$~meV.

We denote this displayed quantity as

\[
R_{\mathrm{bin}}(E',k_\parallel)
=
\min_{E\in\mathrm{bin}(E')}R_{TE}(E,k_\parallel),
\]

where $E'$ is the center of the corresponding energy bin. This binning is used only for visualization and does not change the calculation of the structure itself.

\reportfigure{04_bare_cavity_reflectivity_map.png}{Numerically calculated minimum TE reflectivity inside energy bins of width $0.050$~meV for the microcavity with the exciton resonance switched off. The narrow resonance line corresponds to the photon mode of the finite DBR structure.}{fig:bare-reflectivity}

Figure~\ref{fig:bare-reflectivity} shows one photon branch. Its smooth increase with $k_\parallel$ is the usual behavior of a planar-cavity mode, but here the branch is obtained from the full multilayer calculation rather than from an analytical formula.

For a quantitative comparison, the photon energy in the finite DBR structure is obtained from a photonic eigenmode of the open cavity, not from the ideal-cavity expression. With the chosen transfer-matrix convention, it satisfies

\[
M_{22}(\mathcal E_C^{bg},k_\parallel)=0.
\]

The numerical photon dispersion is given by the real part of the pole, $\operatorname{Re}\mathcal E_C^{bg}(k_\parallel)$. Comparison with the ideal cavity shows how finite DBRs and their reflection phase change the shape of the photon branch.

\reportfigure{05_bare_photon_dispersion_comparison.png}{Photon dispersion: analytical model of an ideal planar cavity and numerical calculation of the finite DBR structure. The dashed curve is obtained analytically from $E_C(k_\parallel)$; the solid curve is the real part of the complex pole $\mathcal E_C^{bg}$ found by the transfer-matrix method from $M_{22}=0$. Energy is shown relative to $E_X(0)$.}{fig:bare-dispersion-comparison}

In Figure~\ref{fig:bare-dispersion-comparison}, the two models almost coincide at small $k_\parallel$, but the numerical branch becomes noticeably steeper at larger $k_\parallel$. This shows the effect of finite DBRs: the reflection phase and the field penetration into the mirrors change with energy and angle, while the ideal-cavity formula does not include these effects.

\clearpage
\section{Quantum well in the microcavity}

The quantum well $W$ has a finite thickness $d_{QW}=6$~nm and refractive index $n_W$. The InGaAs layer itself affects light propagation as an ordinary dielectric. In addition, it contains an exciton transition: near the exciton energy $E_X$, the quantum well interacts strongly with light. The resonant response is described using the single-QW model from Sections~4.4.2 and 4.4.4 of \textit{Microcavities}.

The exciton energy at a given $k_\parallel$ is again

\[
E_X(k_\parallel)
=
E_X(0)
+
\frac{\hbar^2k_\parallel^2}{2M_X},
\]

where $E_X(0)$ is the exciton resonance energy at $k_\parallel=0$, and $M_X$ is the exciton center-of-mass mass in the QW plane.

For the InGaAs layer, the longitudinal component of the wave vector is

\[
q_W(E,k_\parallel)
=
\sqrt{
\left(\frac{n_WE}{\hbar c}\right)^2-k_\parallel^2
}.
\]

We denote the propagation angle inside the QW by $\phi_W$. Then

\[
\cos\phi_W
=
\frac{q_W}{n_WE/(\hbar c)}.
\]

The reflection amplitude of a single QW has a resonant form; see \textit{Microcavities}, Section~4.4.2. In energy units we use

\[
r_X(E,k_\parallel)
=
\frac{i\Gamma_0^{pol}}
{E_X(k_\parallel)-E-i\left(\Gamma_0^{pol}+\gamma_X\right)}.
\]

Here $r_X$ is the amplitude of the reflected wave produced by the exciton resonance. The parameter $\Gamma_0^{pol}$ gives the radiative broadening for the selected polarization, while $\gamma_X$ is an additional homogeneous, non-radiative exciton broadening caused by processes in the material. Both quantities are written as HWHM (half width at half maximum), that is, as half of the full linewidth at half maximum. Therefore, the total resonance HWHM is

\[
\gamma_{QW}^{pol}
=
\Gamma_0^{pol}+\gamma_X,
\]

and the full width at half maximum, FWHM, is

\[
\Gamma_{QW}^{pol}=2\gamma_{QW}^{pol}.
\]

For the physical discussion below, it is enough to remember that increasing $\Gamma_0^{pol}$ or $\gamma_X$ makes the exciton resonance broader.

The transmission amplitude is related to the reflection amplitude by

\[
t_X=1+r_X,
\]

where $t_X$ is the transmission amplitude of the resonant QW.

For oblique propagation, the radiative parameter depends on polarization. According to Section~4.4.4 of the book,

\[
\Gamma_0^{TE}
=
\frac{\Gamma_0}{\cos\phi_W},
\qquad
\Gamma_0^{TM}
=
\Gamma_0\cos\phi_W,
\]

where $\Gamma_0$ is the radiative HWHM at normal incidence. At $k_\parallel=0$, we have $\phi_W=0$, and therefore

\[
\Gamma_0^{TE}=\Gamma_0^{TM}=\Gamma_0.
\]

The exciton resonant response is included in the same transfer-matrix calculation through

\[
T_X
=
\frac{1}{t_X}
\begin{pmatrix}
t_X^2-r_X^2 & r_X\\
-r_X & 1
\end{pmatrix}.
\]

Here $T_X$ describes only the exciton resonance. Propagation through the finite InGaAs thickness is described by

\[
P_W(\ell)
=
\begin{pmatrix}
e^{iq_W\ell} & 0\\
0 & e^{-iq_W\ell}
\end{pmatrix},
\]

where $\ell$ is the propagation distance inside the QW. If the exciton sheet is placed at the center of a layer of thickness $d_{QW}$, the full quantum-well block is

\[
M_{QW}
=
I_{W\rightarrow C}
P_W(d_{QW}/2)
T_X
P_W(d_{QW}/2)
I_{C\rightarrow W}.
\]

The matrices $I_{C\rightarrow W}$ and $I_{W\rightarrow C}$ describe the GaAs--InGaAs and InGaAs--GaAs interfaces. If the exciton resonance is switched off, then

\[
T_X=I,
\]

but the InGaAs layer with refractive index $n_W$ and thickness $d_{QW}$ remains in the structure. This structure defines the background photon resonance $E_C^{bg}$. The index $bg$ means ``background'': the ordinary refractive index of the QW is included, but its exciton resonance is switched off.

The quantum well replaces a GaAs region of the same thickness inside the cavity. Therefore, the total cavity thickness remains constant:

\[
L_C=d_{C_1}+d_{QW}+d_{C_2}.
\]

If $z_{QW}$ is the coordinate of the QW center measured from the left cavity boundary, then

\[
d_{C_1}=z_{QW}-\frac{d_{QW}}{2},
\]

\[
d_{C_2}=L_C-z_{QW}-\frac{d_{QW}}{2}.
\]

Thus, changing the QW position changes $d_{C_1}$ and $d_{C_2}$ but does not increase the total thickness $L_C$.

For a qualitative comparison, the same physical idea can be described by a simpler model of two coupled resonances. Including losses, it can be written as

\[
\left(E_X-E-i\gamma_X\right)
\left(E_C-E-i\frac{\Gamma_C}{2}\right)
=g^2,
\]

where $E_C$ is the photon resonance energy without exciton coupling, $\Gamma_C$ is the photon FWHM, $E_X$ is the exciton energy, and $g$ is the exciton--photon coupling energy. In the limit of small losses,

\[
\Omega_R=2g.
\]

This relation shows the physical meaning of the full model: the exciton and photon resonances again form two coupled branches, as in the $2\times2$ matrix. The difference is that the transfer-matrix method also includes the finite DBR mirrors, the QW position inside the standing-wave field, and the polarization.

The parameter $\Gamma_0$ is determined from the coupling strength of one QW to the light field. For a numerical reference, we use the same experiment by Cilibrizzi et al. with a $2\lambda$ GaAs microcavity and In$_{0.08}$Ga$_{0.92}$As quantum wells. In that sample, three pairs of quantum wells are placed close to field antinodes. Each pair contains two QWs, so the number of strongly coupled wells is $3\times2=6$. Two additional QWs in the same sample are placed close to field minima and mainly collect carriers, so they are not included in this number. For the six strongly coupled QWs, the observed total Rabi splitting was about $8$~meV. If the individual QWs experience approximately the same field, the collective coupling increases approximately as $\sqrt{N}$. Therefore, the reference splitting for one QW is

\[
\Omega_R^{(1)}
\approx
\frac{8}{\sqrt{6}}\ \mathrm{meV}
\approx3.27\ \mathrm{meV}.
\]

Here $N=6$ comes directly from the three QW pairs: $N=3\times2$. The value $\Omega_R^{(1)}$ is an estimate of the splitting that one such QW would produce in a comparable field. The parameter $\Gamma_0$ is then adjusted so that at zero detuning, $E_C(0)=E_X(0)$, and at $k_\parallel=0$,

\[
\operatorname{Re}\mathcal E_{UP}(0)
-
\operatorname{Re}\mathcal E_{LP}(0)
=
3.27\ \mathrm{meV},
\]

where $\mathcal E_{LP}$ and $\mathcal E_{UP}$ are the energies of the lower and upper polariton modes including their finite lifetimes. The meaning of this complex-energy notation is introduced below.

\clearpage
\section{Polariton dispersion from the transfer-matrix method}

After the exciton response $T_X$ is switched on, one photon line in the reflectivity spectrum splits into two polariton branches. The lower branch is denoted by $LP$, and the upper branch by $UP$. Their shapes depend on both the photon and exciton dispersions and on light propagation through the finite DBR mirrors.

The same $k_\parallel$ grid and the same original energy grid as above are used for this map. For visualization, a narrower energy bin of $0.025$~meV is used because the two nearby polariton features must now be resolved.

\reportfigure{06_polariton_reflectivity_map_te.png}{Numerically calculated minimum TE reflectivity inside energy bins of width $0.025$~meV after the exciton resonance of the QW is switched on. The two resonant features correspond to the lower and upper polariton branches.}{fig:polariton-reflectivity}

In Figure~\ref{fig:polariton-reflectivity}, the single photon line from Figure~\ref{fig:bare-reflectivity} becomes two branches. The lower branch becomes weaker in reflectivity as $k_\parallel$ increases, which is consistent with Figure~\ref{fig:hopfield}: its photon fraction decreases.

Next, the numerical result is compared with the $2\times2$ model. In this comparison, it is important not to introduce another source of difference through an approximate photon parabola. Therefore, the photon energy in the $2\times2$ matrix is taken from the same finite multilayer structure with the exciton resonance switched off:

\[
E_C(k_\parallel)=\operatorname{Re}\mathcal{E}_C^{bg}(k_\parallel).
\]

The exciton energy is still given analytically by

\[
E_X(k_\parallel)=E_X(0)+\frac{\hbar^2k_\parallel^2}{2M_X}.
\]

The eigenvalues of the $2\times2$ matrix are then compared with the reflectivity minima of the full structure. In this way, the figure mainly tests the photon--exciton coupling model itself: the photon mode in both approaches belongs to the same finite DBR microcavity.

\reportfigure{07_polariton_tmm_vs_2x2.png}{Polariton dispersion: the $2\times2$ model and the numerical transfer-matrix method. The lines are obtained by analytical diagonalization of the $2\times2$ matrix, but $E_C(k_\parallel)$ is taken from the numerically calculated photon mode of the background DBR structure, $\operatorname{Re}\mathcal E_C^{bg}$. The circles show the positions of the $R_{TE}$ minima in the full calculation with the exciton resonance included.}{fig:tmm-vs-2x2}

In Figure~\ref{fig:tmm-vs-2x2}, the points almost lie on the $2\times2$ model curves. This means that, in the selected range, the simple model of two coupled modes describes the exciton--photon interaction well when the photon dispersion is taken from the same finite DBR cavity used in the numerical model.

\clearpage
\section{Cavity field and QW position}

The QW position affects its coupling to the photon: the coupling is stronger where the standing-wave field is stronger. For the cavity without exciton--photon coupling at $k_\parallel=0$, the standing-wave field is reconstructed. If the amplitudes $A_j^+$ and $A_j^-$ are known at the left boundary of layer $j$, then inside the layer

\[
E_j(\xi)
=
A_j^+e^{iq_j\xi}
+
A_j^-e^{-iq_j\xi}.
\]

For comparison, the normalized intensity is used:

\[
I(z)=
\frac{|E(z)|^2}
{\max_{\mathrm{cavity}}|E(z)|^2}.
\]

The field profile is reconstructed numerically with a step $\Delta z=0.5$~nm. A single QW is first placed at the field antinode closest to the center of the cavity and is then displaced from this position.

To separate the effect of QW position from a change in detuning, for each displacement we first use the same structure with the exciton resonance switched off and retune the photon resonance to the same value $E_C^{bg}(0)$. After switching the exciton response back on, the separation between the real parts of the two polariton poles is found:

\[
\Delta E_{\mathrm{pole}}(\Delta z_{QW})
=
\operatorname{Re}
\left[
\mathcal E_{UP}(0)-\mathcal E_{LP}(0)
\right],
\]

where $\Delta z_{QW}$ is the displacement of the QW center from the selected field antinode.

For one QW, there is a simple expectation: the coupling energy is proportional to the electric-field amplitude at the QW position,

\[
g(z_{QW})\propto |E(z_{QW})|.
\]

Since $I(z)\propto|E(z)|^2$, the estimate for $2g$ can be written as

\[
2g_{\mathrm{field}}(\Delta z)=\Omega_1\sqrt{\frac{I(z_0+\Delta z)}{I(z_0)}}.
\]

Here $\Omega_1$ is the splitting for one QW at the selected antinode, and $z_0$ is the position of this antinode. This expression is not inserted into the numerical model. It is calculated separately and compared with $\Delta E_{\mathrm{pole}}$.

\reportfigure{08_qw_position_coupling.png}{QW position and exciton--photon coupling strength. The upper panel shows the numerically reconstructed normalized intensity $I(z)$ of the photon mode; the points mark the antinodes and the red line marks the selected QW position. The lower panel compares the estimate $2g_{\mathrm{field}}\propto\sqrt{I}$ with the separation between the real parts of the polariton poles obtained from the full transfer-matrix calculation.}{fig:qw-position}

In Figure~\ref{fig:qw-position}, the two dependences almost coincide over most of the range. This directly shows that the coupling strength is controlled by the local field amplitude. Very close to field nodes, the real separation of the poles can decrease faster than the simple estimate because losses become more important when the coupling is very weak.

\clearpage
\section{Effect of the number of DBR pairs}

The finite transmission of a DBR causes photon leakage from the cavity. Increasing the number of mirror pairs increases the reflectivity, makes the resonance narrower, and increases the photon-mode lifetime.

For the complex energy of the photon mode, where $E_C$ denotes the resonance position,

\[
\mathcal E_C=E_C-i\frac{\Gamma_C}{2},
\]

the linewidth is

\[
\Gamma_C=-2\operatorname{Im}\mathcal E_C.
\]

The corresponding cavity Q-factor is

\[
Q_C=\frac{E_C}{\Gamma_C}.
\]

Therefore, a smaller linewidth also means a larger Q-factor and a longer photon-mode lifetime.

For comparison, the following mirror configurations are considered:

\[
(N_t,N_b)=(8,8),\ (12,12),\ (16,16),\ (23,26).
\]

For each number of pairs, the cavity is retuned to the same energy at $k_\parallel=0$. Therefore, the change in $\Gamma_C$ is connected specifically with the mirrors.

In this series, only leakage through finite ideal DBRs is included. Other losses are not modeled.

Therefore, the result describes mirror leakage rather than the full experimental linewidth.

\reportfigure{09_dbr_linewidth_sweep.png}{Photon-mode linewidth caused by leakage through finite DBRs as a function of the number of mirror pairs. All points are obtained numerically from the complex photon pole: $\Gamma_C^{\mathrm{mirror}}=-2\operatorname{Im}\mathcal E_C$. The logarithmic scale shows the rapid decrease of losses as the number of pairs increases.}{fig:dbr-linewidth}

Figure~\ref{fig:dbr-linewidth} shows the expected trend: with more DBR pairs, less light leaks through the mirrors and the resonance becomes narrower. For the $(23,26)$ structure, the calculated linewidth is very small because the model describes idealized mirror leakage rather than all losses present in an experimental sample.

\clearpage
\section{TE--TM splitting}

At normal propagation, TE and TM directions are equivalent. Therefore,

\[
E^{TE}(0)=E^{TM}(0).
\]

At $k_\parallel\neq0$, the phase response of the DBRs and the QW response become polarization dependent. Sections~2.7 and 4.4.4 of \textit{Microcavities} show that TE and TM dispersions can have different curvatures.

For the polariton branches, we define

\[
\Delta_{TE-TM}^{LP}(k_\parallel)
=
E_{LP}^{TE}(k_\parallel)-E_{LP}^{TM}(k_\parallel),
\]

\[
\Delta_{TE-TM}^{UP}(k_\parallel)
=
E_{UP}^{TE}(k_\parallel)-E_{UP}^{TM}(k_\parallel).
\]

Both quantities are obtained from the same numerical calculation. The upper panel shows the four branches themselves, while the lower panel shows their small differences, which are almost invisible on the full energy scale.

The complex TE and TM poles are found on the same $k_\parallel=0\ldots4\ \mu\mathrm m^{-1}$ grid as the other dispersion branches. A separate energy grid is not needed for the poles: the complex root is refined directly for every $k_\parallel$ value.

\reportfigure{10_te_tm_polariton_splitting.png}{TE--TM splitting of the polariton modes. The upper panel shows the numerically found branches $E_{LP}^{TE}$, $E_{UP}^{TE}$, $E_{LP}^{TM}$, and $E_{UP}^{TM}$; the lower panel shows the small differences $\Delta_{TE-TM}^{LP}$ and $\Delta_{TE-TM}^{UP}$. All values are obtained from the complex poles of the same multilayer structure.}{fig:te-tm}

On the full scale of the upper panel in Figure~\ref{fig:te-tm}, the TE and TM branches almost coincide. The lower panel shows that a small splitting is still present and grows to tens of $\mu$eV. For the UP branch, the sign of the difference changes. This only means that the polarization with the slightly higher energy changes as $k_\parallel$ increases.

\clearpage
\section{Open microcavity: mode energies and decay}

The DBR mirrors reflect light very well, but not perfectly. Therefore, photons and polaritons can gradually escape from the microcavity. A resonance has not only a central energy but also a finite linewidth: the faster energy leaves the cavity, the broader the spectral line and the shorter the mode lifetime.

These two quantities can be combined in a complex energy,

\[
\mathcal E(k_\parallel)
=
E(k_\parallel)-i\frac{\Gamma(k_\parallel)}{2}.
\]

Here $E=\operatorname{Re}\mathcal E$ is the resonance energy, and $\Gamma=-2\operatorname{Im}\mathcal E$ is the full width at half maximum. The corresponding characteristic lifetime is

\[
\tau(k_\parallel)=\frac{\hbar}{\Gamma(k_\parallel)}.
\]

To find such a mode, we remove the externally incident wave and look for a field that can exist in the cavity by itself and only leak outward through the mirrors. With the chosen transfer-matrix convention, this condition is

\[
M_{22}(\mathcal E,k_\parallel)=0.
\]

The complex poles are calculated on the same $k_\parallel$ grid with a step of about $0.0167\ \mu\mathrm m^{-1}$. Unlike the reflectivity map, no separate energy grid is scanned here: for each $k_\parallel$, the complex root is refined directly. This description of open modes is discussed in Kavokin et al., \textit{Microcavities}, Section~2.6.

For the lower and upper polaritons, two complex energies $\mathcal E_{LP}$ and $\mathcal E_{UP}$ are obtained. Their real parts give the already familiar branches $E_{LP}$ and $E_{UP}$, while their imaginary parts give the linewidths

\[
\Gamma_{LP,UP}(k_\parallel)
=-2\operatorname{Im}\mathcal E_{LP,UP}(k_\parallel).
\]

Thus, along the same dispersion branch, one can describe both how the polariton energy changes and how quickly the mode decays. A separate figure with the real parts of the complex energies is not needed here because the same dispersion branches have already been shown above. In this section, the main interest is the imaginary part, that is, the linewidth.

For comparison with the simple $2\times2$ model, losses can be added directly to its matrix:

\[
H=
\begin{pmatrix}
E_C-i\Gamma_C/2 & \Omega_R/2\\
\Omega_R/2 & E_X-i\Gamma_X^{nr}/2
\end{pmatrix}.
\]

Here $\Gamma_C$ is the photon-resonance linewidth, and $\Gamma_X^{nr}=2\gamma_X$ is the non-radiative linewidth of the exciton transition. This matrix gives a simple physical interpretation: a polariton that is more photon-like inherits more of the photon losses, while a more exciton-like mode inherits more of the exciton broadening.

In the weak-loss limit, this can be approximately written using the Hopfield coefficients:

\[
\Gamma_{LP}
\approx
|C_{LP}|^2\Gamma_C
+
|X_{LP}|^2\Gamma_X^{nr},
\]

\[
\Gamma_{UP}
\approx
|C_{UP}|^2\Gamma_C
+
|X_{UP}|^2\Gamma_X^{nr}.
\]

These expressions are used only for physical interpretation. The linewidth curves below are not generated from this approximation; they are obtained directly from the imaginary parts of the complex poles of the full multilayer structure.

\reportfigure{11_polariton_linewidths.png}{Linewidths of the lower and upper polaritons. Both curves are obtained numerically from the complex poles of the full structure, $\Gamma_{LP,UP}=-2\operatorname{Im}\mathcal E_{LP,UP}$. The estimate based on the Hopfield coefficients is used only for physical interpretation.}{fig:polariton-linewidths}

The connection with Figure~\ref{fig:hopfield} is especially clear here. At $k_\parallel=0$, both branches have similar photon/exciton composition and similar linewidths. Then the LP becomes more exciton-like and its linewidth approaches the exciton linewidth, while the UP becomes more photon-like and its linewidth decreases toward the narrow photon-mode linewidth shown in Figure~\ref{fig:dbr-linewidth}.

Reflectivity minima and complex poles describe the same open microcavity, but they are extracted from the calculation in different ways. Therefore, we separately check how closely the energy of a reflectivity dip follows the real part of the corresponding pole.

\reportfigure{12_reflectivity_dips_vs_poles.png}{Comparison of reflectivity minima with the real parts of the complex poles. The circles are obtained from minima of the numerically calculated $R_{TE}(E,k_\parallel)$, while the lines come from solving $M_{22}(\mathcal E,k_\parallel)=0$ for the same multilayer structure.}{fig:dips-vs-poles}

Figure~\ref{fig:dips-vs-poles} shows that the reflectivity dips closely follow the real parts of the polariton poles. This is why the points and lines in Figure~\ref{fig:tmm-vs-2x2} also agree well.

\clearpage
\section{Collective coupling of several quantum wells}

Up to this point, one active QW has been considered. In the experimental sample of Cilibrizzi et al., three pairs of quantum wells were placed near antinodes of the standing-wave field. Therefore, six QWs were strongly coupled to the light field. Here, this case is used only as an additional test of collective coupling; the exact epitaxial structure of each QW pair is not reproduced.

First, consider the simple physical idea. If one QW is coupled to the photon mode with coupling energy $g_1$, then for several identical QWs coupled to the same photon mode, the effective coupling is

\[
g_N=
\sqrt{g_1^2+g_2^2+\ldots+g_N^2}.
\]

If all wells are in regions with similar field amplitude, then $g_i\approx g_1$, and therefore

\[
g_N\approx g_1\sqrt{N_{QW}},
\]

while the Rabi splitting should scale as

\[
\Omega_R(N_{QW})\approx\Omega_R(1)\sqrt{N_{QW}}.
\]

This is the analytical expectation. It explains why several QWs increase the coupling: they all interact collectively with the same photon mode.

The numerical model uses $N_{QW}=1,2,4,6$ equivalent active QWs. They are placed in strong-field regions identified from the profile in Figure~\ref{fig:qw-position}. The exact barrier structure inside each experimental QW pair is not specified here: two nearby wells around one antinode are treated as two identical exciton contributions in the same field region. The total cavity thickness is kept constant.

For each $N_{QW}$, the exciton response is first switched off and the photon resonance is retuned to zero detuning. Then the exciton response is switched on, and the two complex polariton poles are found at $k_\parallel=0$. The numerical splitting is defined as

\[
\Omega_R^{\mathrm{num}}(N_{QW})
=
\operatorname{Re}
\left[
\mathcal E_{UP}(0)-\mathcal E_{LP}(0)
\right].
\]

No separate energy grid is used here: the two poles are found directly for each of the four values of $N_{QW}$.

\reportfigure{13_multiple_qw_scaling.png}{Collective enhancement of the coupling for several quantum wells. The circles show $\Omega_R^{\mathrm{num}}$, obtained numerically from the poles of the multilayer structure for $N_{QW}=1,2,4,6$. The line shows the analytical dependence $\Omega_R(1)\sqrt{N_{QW}}$.}{fig:multi-qw}

Figure~\ref{fig:multi-qw} shows that the numerical calculation follows the $\sqrt{N_{QW}}$ dependence. This means that the model reproduces the expected collective enhancement of the coupling.

It is important that the value for one QW was chosen earlier from the experimental reference $8/\sqrt{6}\approx3.27$~meV. Therefore, the point near $8$~meV at $N_{QW}=6$ is not an independent prediction of the experimental result. Here, the test concerns the scaling of the coupling as the number of active QWs increases.

\clearpage
\section{Conclusion}

The analytical picture is compared step by step with the model of the full multilayer microcavity. The $2\times2$ model explains the anticrossing and the LP/UP composition (Figures~\ref{fig:analytical-dispersion} and~\ref{fig:hopfield}). The transfer-matrix method shows the same effects in the finite DBR structure.

The comparison in Figure~\ref{fig:bare-dispersion-comparison} shows that finite mirrors noticeably change the photon dispersion at large $k_\parallel$. At the same time, after this numerical photon branch is used in the $2\times2$ model, the simple model almost coincides with the full polariton calculation (Figure~\ref{fig:tmm-vs-2x2}). This means that the main physics of strong coupling is well described by two coupled modes.

Figure~\ref{fig:qw-position} shows the direct connection between the standing-wave field amplitude and the coupling strength to the QW. Figure~\ref{fig:polariton-linewidths} connects the change of polariton linewidths with the change of their photon and exciton fractions. For several QWs, the numerical calculation reproduces the expected collective scaling $\Omega_R\propto\sqrt{N_{QW}}$ (Figure~\ref{fig:multi-qw}).

Thus, the analytical expressions provide a physical interpretation, while the numerical model includes finite DBRs, QW position, polarization, and the open nature of the microcavity.

\end{document}
